\documentclass[11pt,a4paper]{article}

\usepackage[T1]{fontenc}
\usepackage[utf8]{inputenc}
\usepackage{amsmath}
\usepackage{amssymb}
\usepackage{amsfonts}
\usepackage{mathtools}
\usepackage{xcolor}
\usepackage{hyperref}
\usepackage{listings}
\usepackage{xspace}
\usepackage{geometry}
\usepackage{natbib}
\usepackage{graphicx}

\geometry{margin=1in}

% Lean code formatting
\lstdefinelanguage{Lean}{
  keywords={def, theorem, lemma, structure, class, instance, import, namespace, end, by, sorry, simp, rfl, exact, intro, intros, constructor, cases, induction, have, let, where, if, then, else, match, with},
  sensitive=true,
  comment=[l]{--},
  morecomment=[s]{/-}{-/},
  morestring=[b]",
  basicstyle=\ttfamily\small,
  keywordstyle=\color{blue},
  commentstyle=\color{gray},
  stringstyle=\color{red},
  breaklines=true,
  showstringspaces=false,
  tabsize=2
}

\lstset{language=Lean, basicstyle=\ttfamily\small, extendedchars=true,
  literate=
    {→}{{$\rightarrow$}}1
    {∀}{{$\forall$}}1
    {∃}{{$\exists$}}1
    {≠}{{$\neq$}}1
    {≃}{{$\simeq$}}1
    {≅}{{$\cong$}}1
    {∪}{{$\cup$}}1
    {∩}{{$\cap$}}1
    {×}{{$\times$}}1
    {∈}{{$\in$}}1
    {Φ}{{$\Phi$}}1
    {τ}{{$\tau$}}1
    {λ}{{$\lambda$}}1
    {∘}{{$\circ$}}1
    {₀}{{\textsubscript{0}}}1
    {₁}{{\textsubscript{1}}}1
    {₂}{{\textsubscript{2}}}1
    {₃}{{\textsubscript{3}}}1
}

% Notation
\newcommand{\Future}{\mathcal{F}}
\newcommand{\seqBind}{\mathbin{>>=}}
\newcommand{\parTensor}{\otimes}
\newcommand{\fork}{\mathbin{|}}
\newcommand{\merge}{\oplus}
\newcommand{\Aff}{\Phi}
\newcommand{\State}{\mathcal{S}}
\newcommand{\Traj}{\tau}
\newcommand{\id}{\text{id}}
\newcommand{\well}{\text{wf}}

\title{Composable Future: An Algebraic Theory of Paradigmatic Transitions}
\author{I Made Agus Kresna Sucandra\\
Fakultas Kedokteran, Universitas Udayana}
\date{April 2026}

\begin{document}

\maketitle

\begin{abstract}
We propose \emph{Composable Future}, a formal theory of paradigmatic transitions as composable algebraic structures. A future is a 4-tuple $F = (S_0, \tau, S_1, \Phi)$ where $S_0$ and $S_1$ are paradigmatic states, $\tau$ is a trajectory of change, and $\Phi$ is an affordance set of accessible futures from $S_1$. We define four primitive operators (sequential bind, parallel tensor, fork, merge) and prove fundamental laws: identity (substantive in the affordance component, ADR-0005), closure, well-formedness preservation, and associativity. The trajectory carries an explicit \emph{path} of intermediate states; sequential bind concatenates paths, so associativity is \emph{substantive} --- it holds by \texttt{List.append\_assoc} on the concatenated paths, not by definitional collapse of trajectories --- and is further structured via an indexed monad construction for path-dependent trajectories. We resolve the central open problem (associativity under path-dependence) by carrying path-dependence in the trajectory itself, complemented by Orchard et al.'s indexed monad framework, and establish the correct equivalence relation between futures (OP3): path-trace isomorphism implemented as \texttt{FutureIso} + \texttt{PathIso} + \texttt{TrajectoryEquiv} in Lean~4, giving \texttt{ComposableFuture} the structure of a setoid. The theory occupies the intersection of category theory, paradigm change studies, process algebra, affordance theory, and futures formalization — a structure not currently named or investigated in any of these domains.
\end{abstract}

\section{Introduction}

The study of paradigmatic transitions — moments when entire epistemic structures, constraint systems, and research infrastructures change — remains largely qualitative. Kuhn's \emph{Structure of Scientific Revolutions} \cite{kuhn1962} established the paradigm as a unit of analysis, and Lakatos \cite{lakatos1978} formalized the distinction between progressive and degenerative research programmes. Yet neither framework provides operators for composing paradigms, laws for their interaction, or a formal account of what futures become possible after a transition.

We propose that paradigmatic futures have the algebraic properties of composable types. They can be combined without either being destroyed, sequenced without loss of identity, and the result of composition is itself a valid future that can be further composed. This is distinct from:
\begin{itemize}
\item \emph{Convergence} — which implies two things merging into one fixed outcome
\item \emph{Futures studies} — which models scenarios qualitatively, not algebraically
\item \texttt{Future<T>} in async programming — which operates on computations, not paradigms
\end{itemize}

\subsection{The Core Claim}

A \emph{composable future} is a 4-tuple:
\[
F = (S_0, \tau, S_1, \Phi)
\]
where:
\begin{itemize}
\item $S_0$ is the current paradigmatic state (existing assumptions, constraints, infrastructure)
\item $\tau$ is a trajectory — the mechanism of change from $S_0$ to $S_1$
\item $S_1$ is a reachable paradigmatic state
\item $\Phi$ is an affordance set — what new compositions become possible from $S_1$
\end{itemize}

The central contribution is proving that these futures compose associatively, and that the composition laws (identity, closure, associativity, well-formedness preservation) hold formally. Path-dependence --- the fact that $\tau$ may depend on the history of prior transitions --- is carried explicitly in the trajectory's \texttt{path} field; sequential bind concatenates paths, so associativity is a substantive consequence of list-append associativity (not a definitional artifact), and is additionally structured via an indexed monad construction.

\subsection{Contributions}

\begin{enumerate}
\item \textbf{Formal definition} of composable futures as 4-tuples with well-formedness conditions
\item \textbf{Four primitive operators} (sequential bind, parallel tensor, fork, merge) with categorical semantics
\item \textbf{Proof of associativity} for all futures (general case) and for stateless futures (restricted case)
\item \textbf{Indexed monad construction} resolving associativity under path-dependence via Orchard et al.'s framework
\item \textbf{Weak associativity theorems} at the affordance level, holding even when trajectories differ
\item \textbf{Equivalence relation for futures} (OP3): \texttt{FutureIso} pairs state bijections with path-trace trajectory equivalence (\texttt{PathIso} + \texttt{TrajectoryEquiv}), forming an equivalence relation and setoid. \texttt{parTensor} is commutative up to this isomorphism via the symmetric-monoidal braiding.
\item \textbf{Open problems inventory} identifying 5 gaps in existing literature and 12 internal research directions
\end{enumerate}

\section{Related Work}

The theory of composable futures draws from five adjacent domains, none of which currently addresses the intersection.

\subsection{Category Theory Applied to Complex Systems}

Applied category theory provides compositional machinery for open systems. Spivak \cite{spivak2014} and Fong \& Spivak \cite{fong2018} establish category theory as a cross-domain modeling framework. Furter et al. \cite{furter2025} extend symmetric monoidal categories to handle uncertainty via Markov kernels, treating probabilistic transitions as composable structures. However, none of these works apply categorical composition to paradigmatic transitions — the level at which entire epistemic structures change.

\subsection{Formal Models of Paradigm Change}

Kuhn \cite{kuhn1962} and Lakatos \cite{lakatos1978} provide qualitative frameworks for paradigm change. Bechberger \& Kühnberger \cite{bechberger2017} formalize paradigmatic states geometrically using conceptual spaces. Yet these frameworks provide no operators, no composition laws, and no formal account of affordances over states.

\subsection{Process Algebra and Concurrent Systems}

Hoare \cite{hoare1985}, Milner \cite{milner1989}, and Katis et al. \cite{katis2009} develop mature operator vocabularies (sequential bind, parallel tensor, fork, merge) with categorical semantics. However, these operators apply exclusively to computational processes, not paradigmatic ones.

\subsection{Affordance Theory}

Chemero \cite{chemero2003}, Turvey \cite{turvey1992}, Şahin et al. \cite{sahin2007}, and Kyriakoullis \& Zaphiris \cite{kyriakoullis2016} establish the relational and dispositional ontology required for affordance sets. Yet affordances are not formalized as typed functions over abstract state transitions.

\subsection{Futures Formalization}

Iacona \& Iaquinto \cite{iacona2021} formalize what it means to \emph{believe} a future in branching-time logic. However, they do not address what it means for futures to \emph{compose}.

\subsection{Gap Statement}

Applied category theory provides compositional machinery for open systems, and indexed monad theory provides path-dependent composition semantics, but neither framework has been applied to paradigmatic transitions. The formal treatment of paradigm change formalizes states geometrically or classifies transitions qualitatively, but provides no operators or composition laws. Process algebra provides a mature operator vocabulary with categorical semantics, but applies these operators exclusively to computational processes. Affordance theory establishes the ontology required for affordance sets, but does not define them as typed functions over abstract state transitions. The formal treatment of future contingents formalizes what it means to believe a future, but not what it means for futures to compose.

\emph{Composable Future occupies the intersection of all five domains — a structure not currently named, formalized, or investigated in any of them.}

\section{Core Definitions}

\subsection{Paradigmatic State}

A paradigmatic state $S$ is a triple of types:
\[
S = (\text{assumptions}, \text{constraints}, \text{infrastructure})
\]
Each component is a type (in the sense of dependent type theory), not a concrete value. This allows $S$ to be defined at the specification level, before any particular instantiation.

\textbf{Remark (State Identity Criterion).} Two paradigmatic states are
\emph{equal} when their three component types are propositionally equal:
$S = S' \iff S.\text{assumptions} = S'.\text{assumptions} \land
S.\text{constraints} = S'.\text{constraints} \land
S.\text{infrastructure} = S'.\text{infrastructure}$. This strict criterion is
deliberately fine: $A \times B$ and $B \times A$ are \emph{not} equal under it
(no univalence). The intended working notion of ``same paradigm'' is the
\emph{weaker} component-wise bijection $\simeq$, captured by \texttt{StateIso}
(and lifted to futures by \texttt{FutureIso}), under which \texttt{ComposableFuture}
forms a setoid (\S\ref{sec:op3}). Equality is for construction; isomorphism is
for identity of paradigms.

\subsection{Trajectory}

A trajectory $\tau$ is a transition between paradigmatic states, carrying an
internal path of intermediate states (ADR-0002):
\[
\tau = (\text{source} : S_0,\ \text{path} : \text{List}~\mathcal{S},\ \text{target} : S_1)
\]
The trajectory encodes the mechanism of change. In the stateless case the path
is empty and $\tau$ is path-independent; in the general case the \texttt{path}
field records the sequence of intermediate paradigmatic states, making
path-dependence a first-class part of the trajectory rather than an external
effect index. Sequential composition concatenates paths (\S\ref{sec:assoc}),
which is what makes associativity \emph{substantive} rather than definitional.

\subsection{Affordance Set}

An affordance set specifies what compositions become possible from a state.
Formally it is a map into the powerset of futures:
\[
\Phi : S_1 \to \mathcal{P}(\Future)
\]
where $\Phi(S) = \texttt{AffordanceSet}~S = \{F \in \Future : F.S_0 = S\}$ is
the set of all futures whose source is $S$. This is well-defined before $S_1$
is concretely realized (OP2, resolved) and closed under composition (OP4,
resolved).

\textbf{Remark (Affordance Circularity and Its Lean Realization).} The
definition is prima facie circular: a future $F$ contains $\Phi$, and
$\Phi : S_1 \to \mathcal{P}(\Future)$ ranges over futures. In set theory this
is handled by stratification; in a proof assistant it is a strict-positivity
question. A naively \emph{stored} field $\Phi : \texttt{Set}~\Future$ is
\textbf{not} admissible in Lean~4: $\texttt{Set}~T = T \to \texttt{Prop}$
places $\Future$ in a negative (contravariant) position --- a strict-positivity
violation that the Lean kernel rejects (verified 2026-05-15). The
formalization therefore stores the affordance set \emph{state-anchored} as
$\Phi : \texttt{Set}~\mathcal{S}$ (a set of \emph{anchor states}) and recovers
the powerset-of-futures view by the projection
$\texttt{afforded}~F := \{G : G.S_0 \in F.\Phi\}$. For a well-formed future
this is provably content-equivalent to the intended object,
$\texttt{afforded}~F = \texttt{AffordanceSet}~F.S_1$
\big[Lean: \texttt{afforded\_eq\_affordanceSet}\big]. The literal
$\texttt{Set}~\Future$ field is recoverable only via a coinductive
\texttt{ComposableFuture}, which is deferred indefinitely; the state-anchored
encoding is the faithful, kernel-accepted realization (ADR-0005). This
resolves the circularity concern raised against the v0.1 preprint.

\subsection{Composable Future}

A composable future is a 4-tuple:
\[
F = (S_0, \tau, S_1, \Phi)
\]
with a well-formedness condition (ADR-0005):
\[
\well(F) \equiv \tau.\text{source} = S_0 \ \land\ \tau.\text{target} = S_1
            \ \land\ F.\Phi = \texttt{AffordanceSet}~S_1
\]
The third conjunct anchors the affordance set to the realized target state:
a well-formed future affords exactly the futures accessible from $S_1$. It is
this conjunct that makes the identity laws (\S\ref{sec:laws}) substantive
rather than definitional --- the Lean proof of right identity uses it
essentially (it is not discharged by \texttt{rfl} or a \texttt{Subsingleton}
instance).

\subsection{Interpretations of $F$}

The 4-tuple admits three compatible readings. They are instantiations of one
structure, not competing definitions:
\begin{itemize}
\item \textbf{Morphism reading} (\S\S3--4): $F$ is a categorical morphism
  $S_0 \to S_1$; $\seqBind$ is composition; the laws are the category axioms.
\item \textbf{Affordance reading} (\S\ref{sec:op3} and the affordance set): $\Phi$
  is a relational structure over states (Chemero \cite{chemero2003}); identity
  of paradigms is \texttt{FutureIso}, not strict equality.
\item \textbf{Probabilistic reading} (\S6): $\tau$ is a Markov kernel
  $S_0 \to \mathcal{P}(S_1)$ and composition is Kleisli bind (Furter et al.
  \cite{furter2025}).
\end{itemize}
Each subsequent section fixes one reading; the formalization shows them
mutually consistent (the same operators and laws hold under all three).

\section{Operators and Laws}

\subsection{Primitive Operators}

\subsubsection{Sequential Bind}
\[
F \seqBind G \quad \text{(when } F.S_1 = G.S_0 \text{)}
\]
Constructs a new future by composing $F$ and $G$ sequentially, concatenating
trajectory paths:
\[
F \seqBind G = \big(F.S_0,\ \{F.\tau.\text{source},\
F.\tau.\text{path} \mathbin{+\!\!+} G.\tau.\text{path},\
G.\tau.\text{target}\},\ G.S_1,\ G.\Phi\big)
\]

\subsubsection{Parallel Tensor}
\[
F \parTensor G
\]
Both proceed in the joint paradigm: states combine by component-wise
cartesian product, and the affordance set is the product
$\Phi^{F} \times \Phi^{G} = \{\,\texttt{paradigmaticTensor}~a~b : a \in F.\Phi,\
b \in G.\Phi\,\}$. (Open Problem~2, state product structure, is
\textbf{resolved}; commutativity is up to the symmetric-monoidal braiding,
\S\ref{sec:op3}.)

\subsubsection{Fork}
\[
F \fork G
\]
A branch point where one of two futures is realized; the result carries
$\Phi^{F} \cup \Phi^{G}$ (coproduct of affordances). \textbf{Remark
(Deferral).} Paper~1 fixes the left-biased reading sufficient for the
algebra. \emph{When} branching occurs, whether it is observer-relative, and
whether the unrealized branch's history coexists are parameters of the
theory; their formal account requires the time enrichment of Paper~2 and is
deliberately deferred.

\subsubsection{Merge}
\[
F \merge G
\]
Two independent futures reconverge; the result carries $\Phi^{F} \cap \Phi^{G}$
(symmetric convergence). \textbf{Remark (Scope).} The present definition covers
the \emph{symmetric} case only. \emph{Absorptive} merge --- asymmetric
resource transfer with termination of the absorbed source --- is a Paper~2
question, requiring the resource enrichment of Coecke, Fritz, and Spekkens
\cite{coecke2016}.

\subsubsection{Identity}
\[
\id_S
\]
The null future — a transition that changes nothing:
\[
\id_S = (S,\ \{\text{source} := S,\ \text{path} := [\,],\ \text{target} := S\},\ S,\ \Phi(S))
\]
where $\Phi(S) = \texttt{AffordanceSet}~S$ is the set of all futures accessible
from $S$ (\textbf{Option B}, ADR-0005). A transition that changes nothing
changes nothing about what is accessible, so the null future \emph{preserves}
the full affordance set rather than zeroing it.

\textbf{Remark 4.1 (Null Future Preserves Affordances; revising the v0.1
preprint).} The earlier preprint stated that the affordance set of
$F \seqBind \id_{S_1}$ is $\Phi_\emptyset = \emptyset$. This is corrected:
\[
\Phi(F \seqBind \id_{S_1}) = F.\Phi,
\]
because composing with the null future preserves all affordances accessible
from $S_1$ --- a transition that changes nothing changes nothing about what is
accessible (Chemero's relational ontology of affordances \cite{chemero2003}
supports this: if neither the composition context nor the state features
change, the relational structure is preserved). The operation that genuinely
zeros affordances is the \emph{terminate} operator, a fifth unary operator
deferred to Paper~2, where it becomes substantive under resource enrichment
(Coecke, Fritz, and Spekkens \cite{coecke2016}). The v0.1 Remark~4.1 conflated
identity with termination; Option~B separates them.

\subsection{Fundamental Laws}\label{sec:laws}

\subsubsection{Left Identity}
\[
\id_{S_0} \seqBind F = F
\]
\textbf{Proved} for well-formed $F$ \big[Lean: \texttt{Laws.left\_identity}\big].
Formerly stated as Open Problem~9 (pending affordance structure); resolved by
ADR-0005. With Option~B, $\id_{S_0}$ carries $\Phi = \texttt{AffordanceSet}~S_0$
and $\seqBind$ propagates the right operand's affordances, so both sides carry
$F.\Phi$; the path equality is $[\,]\mathbin{+\!\!+}F.\tau.\text{path} =
F.\tau.\text{path}$. No \texttt{Subsingleton} guard is required.

\subsubsection{Right Identity}
\[
F \seqBind \id_{S_1} = F
\]
\textbf{Proved} for well-formed $F$ \big[Lean: \texttt{Laws.right\_identity}\big].
Formerly Open Problem~10; resolved by ADR-0005. The proof is \emph{substantive}
in the affordance conjunct of $\well(F)$: $F \seqBind \id_{S_1}$ carries
$\Phi = \id_{S_1}.\Phi = \texttt{AffordanceSet}~S_1$, and matching this against
$F.\Phi$ uses $\well(F)$'s third conjunct $F.\Phi = \texttt{AffordanceSet}~S_1$
essentially (verified by \texttt{\#print}: the proof term depends on this
hypothesis, not on \texttt{rfl} or \texttt{Subsingleton}; it uses only the
\texttt{propext} axiom).

\subsubsection{Closure}
\[
\forall F, G \in \Future : F \seqBind G \in \Future
\]
\textbf{Proved} \big[Lean: \texttt{Laws.closure}\big]. Any composition of valid futures is itself a valid future.

\subsubsection{Well-Formedness Preservation}
\[
\well(F) \land \well(G) \land F.S_1 = G.S_0 \implies \well(F \seqBind G)
\]
\textbf{Proved} \big[Lean: \texttt{Laws.seqBind\_well\_formed}\big]. Sequential
bind preserves all three well-formedness conjuncts: the source comes from $F$,
the target and the affordance anchor ($\Phi = G.\Phi = \texttt{AffordanceSet}~G.S_1$)
come from $G$.

\subsubsection{Associativity}
\[
(F \seqBind G) \seqBind H = F \seqBind (G \seqBind H)
\]
\textbf{Proved} \big[Lean: \texttt{Laws.seqBind\_assoc}\big]. Sequential bind
\emph{concatenates} the trajectory paths
($(F \seqBind G).\tau.\text{path} = F.\tau.\text{path} \mathbin{+\!\!+}
G.\tau.\text{path}$), so both groupings produce the future
\[
(S_0,\ \{\text{source} := F.\tau.\text{source},\
\text{path} := F.\tau.\text{path} \mathbin{+\!\!+} G.\tau.\text{path}
\mathbin{+\!\!+} H.\tau.\text{path},\
\text{target} := H.\tau.\text{target}\},\ H.S_1,\ H.\Phi)
\]
and equality follows from associativity of list append. The proof term is
\emph{not} \texttt{rfl}: it is \texttt{List.append\_assoc}.

\section{Associativity: Main Technical Results}\label{sec:assoc}

\subsection{General Associativity}

\textbf{Theorem (General Associativity).} For all composable futures $F, G, H$ with compatible states:
\[
(F \seqBind G) \seqBind H = F \seqBind (G \seqBind H)
\]

\textbf{Proof.} Sequential bind builds the composed trajectory by
concatenating the operands' \texttt{path} fields (ADR-0002). The two groupings
yield paths $(F.\tau.\text{path} \mathbin{+\!\!+} G.\tau.\text{path})
\mathbin{+\!\!+} H.\tau.\text{path}$ and $F.\tau.\text{path} \mathbin{+\!\!+}
(G.\tau.\text{path} \mathbin{+\!\!+} H.\tau.\text{path})$ respectively; these
are equal by \texttt{List.append\_assoc}, and the remaining components
($S_0$, endpoints, $\Phi = H.\Phi$) coincide. The proof in Lean is:
\begin{lstlisting}
theorem seqBind_assoc (F G H : ComposableFuture)
  (h₁ : F.S₁ = G.S₀) (h₂ : G.S₁ = H.S₀) :
  seqBind (seqBind F G h₁) H (by exact h₂) =
  seqBind F (seqBind G H h₂) (by exact h₁) := by
  simp [seqBind, List.append_assoc]
\end{lstlisting}

This is the \emph{substantive} reading of associativity: paradigm
trajectories associate because their recorded histories associate under
concatenation --- not because trajectories are collapsed to endpoints. The
v0.1 preprint claimed a definitional (\texttt{rfl}) proof that discarded
intermediate history; ADR-0002 corrected this by making \texttt{path} a
first-class field, so the result now genuinely covers path-dependent
trajectories rather than sidestepping them.

\subsection{Stateless Case}

\textbf{Definition (Stateless).} A trajectory $\tau$ is stateless if it does
not depend on the history of prior transitions; in this formalization that is
exactly an empty path, $\tau.\text{path} = [\,]$.

\textbf{Theorem (Stateless Associativity).} For stateless futures $F, G, H$:
\[
(F \seqBind G) \seqBind H = F \seqBind (G \seqBind H)
\]

This is the specialization of the general result to empty paths, where
\texttt{List.append\_assoc} degenerates to $[\,] = [\,]$. Strict associativity
thus holds in the path-independent case as the trivial instance of the
substantive theorem.

\subsection{Indexed Monad Construction}

While the general case admits strict associativity, we provide an indexed monad framework for fine-grained tracking of trajectory types and path-dependence.

\subsubsection{Trajectory Type}

A trajectory type $t$ is a classification of trajectory behavior:
\[
t = (\text{name} : \text{String}, \text{isStateless} : \text{Bool})
\]

\subsubsection{Trajectory Type Composition}

We define a monoid structure on trajectory types via a typeclass:
\begin{lstlisting}
class TrajectoryTypeCompose (T : Type) where
  compose : T → T → T
  unit : T
  assoc : ∀ a b c, compose (compose a b) c = 
                    compose a (compose b c)
  left_id : ∀ a, compose unit a = a
  right_id : ∀ a, compose a unit = a
\end{lstlisting}

The composition rule is: stateless acts as the unit; composition of non-stateless types produces a combined type.

\subsubsection{Indexed Future}

An indexed future is graded by trajectory type:
\[
\text{IndexedFuture}(t) = (S_0, S_1, \tau, \Phi, \well)
\]
where the index $t$ tracks the trajectory type.

\subsubsection{Indexed Associativity}

\textbf{Theorem (Indexed Associativity).} For indexed futures $F : \text{IndexedFuture}(t_1)$, $G : \text{IndexedFuture}(t_2)$, $H : \text{IndexedFuture}(t_3)$:
\[
\text{cast}(\text{assoc}(t_1, t_2, t_3)) \left( (F \seqBind G) \seqBind H \right) = F \seqBind (G \seqBind H)
\]

The cast uses the associativity law of trajectory type composition. This provides a structured way to track path-dependence at the type level, even though the underlying futures compose associatively.

\subsection{Weak Associativity}

\textbf{Definition (Future Equivalence).} \texttt{FutureEquiv} is a
deliberately weak, \emph{state-level} equivalence: two futures are equivalent
when their source and target states agree,
\[
F \equiv G \iff F.S_0 = G.S_0 \ \land\ F.S_1 = G.S_1
\]
(trajectories may differ, capturing path-dependence). It deliberately does
\emph{not} assert $\Phi$ equality --- with the stored affordance field
(ADR-0005), same-$S_1$ does not imply same-$\Phi$ in general (only well-formed
futures satisfy $\Phi = \texttt{AffordanceSet}~S_1$). The full
$\Phi$-aware notion is \texttt{FutureIso} (\S\ref{sec:op3}); \texttt{FutureEquiv}
isolates the coarser structural invariant.

\textbf{Theorem (Weak Associativity).} For all composable futures $F, G, H$:
\[
(F \seqBind G) \seqBind H \equiv F \seqBind (G \seqBind H)
\]
\big[Lean: \texttt{WeakAssoc.weak\_assoc\_affordance}\big]. Both groupings share
$S_0 = F.S_0$ and $S_1 = H.S_1$, so they are \texttt{FutureEquiv} regardless of
the differing trajectory paths. (They additionally share $\Phi = H.\Phi$ via
$\seqBind$'s affordance propagation, though \texttt{FutureEquiv} does not
record this.) Even when path-dependence makes the trajectories differ, the
state structure composes associatively.

\section{Probabilistic Extension (Preview)}

Phase 3 extends composable futures to probabilistic trajectories via Kleisli categories. Following Furter et al. \cite{furter2025}, we define:
\[
\tau : S_0 \to \mathcal{P}(S_1)
\]
as a Markov kernel, and composition via Kleisli bind. The Kleisli category is known to be associative, so the probabilistic extension inherits associativity from the monad structure.

\section{Open Problems}

\subsection{OP2: Well-Definedness of $\Phi$ Before $S_1$ is Realized — \textbf{Resolved}}

Turvey \cite{turvey1992} notes that affordances are dispositional --- potential
before actual. \textbf{Resolved} (ADR-0005): $\Phi$ is a stored field whose
value, $\texttt{AffordanceSet}~S = \{F : F.S_0 = S\}$, is a set comprehension
determined by $S$ alone. It is well-defined and non-empty (it always contains
$\id_S$) before any concrete $S_1$ is realized, because it is a specification-level
set, not a value-level computation that depends on $S_1$ being instantiated
\big[Lean: \texttt{affordanceSet\_nonempty}\big].

\subsection{OP3: Correct Equivalence Relation Between Futures — \textbf{Resolved}}\label{sec:op3}

We establish path-trace isomorphism as the correct equivalence. \texttt{StateIso}~$S~T$ records component-wise bijections ($\simeq$) on the three fields of a paradigmatic state. \texttt{PathIso}~$xs~ys$ is an inductive type witnessing a pointwise \texttt{StateIso} along the intermediate state lists of two trajectories — the list-indexed analogue of strong bisimulation for deterministic systems. \texttt{TrajectoryEquiv}~$\tau_1~\tau_2$ combines \texttt{StateIso} on source and target with \texttt{PathIso} on the path. \texttt{FutureIso}~$F~G$ pairs state isomorphisms on $S_0$/$S_1$ with a \texttt{TrajectoryEquiv} on $\tau$.

The choice of strong (rather than weak) bisimulation is justified by the absence of silent transitions in \texttt{Trajectory}; Wang's history-preserving bisimilarity is redundant because \texttt{Trajectory} has no branching structure. \texttt{FutureIso} is an equivalence relation (\texttt{refl}, \texttt{symm}, \texttt{trans}) and \texttt{ComposableFuture} forms a \texttt{Setoid} under \texttt{Nonempty~(FutureIso~$\cdot$~$\cdot$)}.

The main result is \texttt{parTensor\_comm\_iso}: $F \otimes G \cong G \otimes F$ via \texttt{Equiv.prodComm} on state components and \texttt{PathIso.nil} on the trajectory (both \texttt{parTensor} trajectories have empty paths). The state- and trajectory-level isomorphism is mechanized in Lean~4 with no new axioms and no \texttt{sorry}. The single exception is the affordance component of this one witness (\texttt{parTensor\_comm\_iso.phi}): strict equality of the two product affordance sets reduces to type-level product commutativity $A \times B = B \times A$, which is independent of Lean's axioms (unprovable without univalence, irrefutable without product-constructor injectivity). It is recorded as a single documented Phase-4 \texttt{sorry}; it is the same limitation as the conditional non-commutativity result for $\otimes$ (strict $A \otimes B \neq B \otimes A$ requires product-constructor injectivity, not available without an added axiom) and is not a gap in the equivalence relation itself. The categorically correct statement --- commutativity \emph{up to} the symmetric-monoidal braiding --- is fully proved by \texttt{parTensor\_comm\_iso} at the state/trajectory level.

\subsection{OP4: Composition of Affordance Sets — \textbf{Resolved}}

Şahin et al. \cite{sahin2007} demonstrate affordance chaining. \textbf{Resolved}
(ADR-0005): sequential composition propagates the right operand's affordances,
$\Phi(F \seqBind G) = G.\Phi$, so the composed state's affordance set is
well-defined and equals the affordances of the final state
\big[Lean: \texttt{seqBind\_$\Phi$\_eq}, \texttt{seqBind\_mem\_affordanceSet}\big].
Affordance chaining is therefore closed under composition.

\subsection{OP5: Completeness}

Are all paradigmatic futures reachable by finite composition from a given
$S_0$? The trivial form is closed by the type system (every well-formed future
with source $S_0$ is, by definition, a member of $\texttt{AffordanceSet}~S_0$).
The non-trivial form --- a characterization of which transitions are
\emph{externally} realizable --- requires a semantic model beyond the
formalization and is deferred (Paper~2 enrichment territory).

\subsection{Internal Open Problems — \textbf{Resolved by ADR-0005}}

\begin{itemize}
\item OP8 (empty affordance set for identity): \textbf{resolved} ---
  Option~B sets $\id_S.\Phi = \texttt{AffordanceSet}~S$ (non-empty); the
  empty-affordance operator is \emph{terminate}, deferred to Paper~2.
\item OP9 (left identity proof): \textbf{proved}
  \big[\texttt{Laws.left\_identity}\big].
\item OP10 (right identity proof): \textbf{proved}
  \big[\texttt{Laws.right\_identity}\big].
\item OP17 (indexed-monad identity affordance): \textbf{resolved} ---
  the indexed-future identity inherits Option~B's affordance assignment.
\end{itemize}

\section{Operational Falsifiability}

The theory yields a concrete, falsifiable test distinguishing a genuine
paradigm shift from a rebranded composition:

\begin{quote}
A claimed ``new paradigm'' is a \emph{composition} if it can be expressed as a
finite term in the pre-existing operator set
$\{\seqBind, \parTensor, \fork, \merge\}$ without loss. It is an
\emph{extension} only if its expression requires a new operator that did not
previously exist.
\end{quote}

This is not merely rhetorical: the operator set is fixed, so membership is
decidable in principle for a given candidate. \textbf{Worked instances.} Genuine
extensions introduced operators absent from the prior vocabulary:
backpropagation (1986), the attention mechanism (2017), and score-matching /
diffusion (2020) each added an irreducible compositional primitive. By
contrast, ``Agentic AI'' is a composition:
\[
\text{Agent} \;=\; \text{LLM} \seqBind \text{tool-call} \seqBind
\text{retrieval} \fork \text{control-flow}
\]
expressible entirely in the existing operators with no loss --- hence an
engineering recombination, not a paradigm shift under this theory. The test
makes the central claim refutable: exhibit a widely-acknowledged paradigm
shift that is expressible as a finite operator term (falsifying ``real shifts
need new operators''), or a rebranding that provably is not (falsifying the
sufficiency of the operator set).

\section{Conclusion}

We have proposed composable futures as a formal theory of paradigmatic
transitions, defined four primitive operators, and proved fundamental laws.
The key insight is that associativity is \emph{substantive}: the trajectory
carries an explicit path of intermediate states, sequential bind concatenates
paths, and associativity follows from \texttt{List.append\_assoc} --- the proof
genuinely covers path-dependent trajectories rather than collapsing them to
endpoints. Identity is likewise substantive in its affordance component
(ADR-0005, Option~B). The theory occupies the intersection of five adjacent
domains and resolves the central open problem (associativity under
path-dependence) by carrying path-dependence in the trajectory itself,
complemented by Orchard et al.'s indexed monad framework.

\textbf{Status.} OP1--OP4 are resolved and mechanized; OP3 (equivalence
relation) is closed via \texttt{FutureIso}. The non-commutativity of $\otimes$
is resolved as a recorded decision: strict $A \otimes B \neq B \otimes A$ is
independent of Lean's axioms (unprovable without univalence, irrefutable
without product-constructor injectivity), and the categorically correct
statement --- commutativity up to the symmetric-monoidal braiding --- is fully
proved (\texttt{parTensor\_comm\_iso}). The Lean development has 0 errors and a
single documented Phase-4 \texttt{sorry} (\texttt{parTensor\_comm\_iso.phi}),
exactly the type-level product-commutativity limitation above; no new axioms.

\textbf{Toward Paper~2 (Enriched Composable Future).} The trajectory will be
enriched with time (Lawvere enrichment over $(\mathbb{R}_{\geq 0}, +, 0)$) and
resources (Coecke, Fritz, and Spekkens \cite{coecke2016}). This makes the
fifth, unary \emph{terminate} operator substantive --- it is the operation
that genuinely zeros affordances, distinct from identity (\textbf{Remark 4.1})
--- and provides the formal account of fork timing and \emph{absorptive}
merge, both deferred from Paper~1's symmetric scope. Paper~3 applies the
enriched machinery to systems-leverage analysis.

The codebase is available at
\url{https://github.com/SHA888/composable-future}, with all proofs verified in
Lean~4 (0 errors; one documented Phase-4 \texttt{sorry} as noted; no new
axioms).

\bibliographystyle{plainnat}
\bibliography{references}

\end{document}
